\documentclass[12pt]{article}
\usepackage{amsmath}
\usepackage{amssymb}
\usepackage[a4paper,margin=1in]{geometry}

\begin{document}

\section*{Lecture 4: Joint Probability and Conditional Probability}

\textbf{Prepared strictly from the provided lecture material} \\
\textbf{Source:} 

L4

This scribe reconstructs the lecture exactly from the provided material. \\
All concepts, definitions, assumptions, axioms, corollaries, and propositions are stated with their logical dependencies exactly as presented. \\
No external examples or explanations are introduced.

\section*{1. Course Context and Motivation}

\subsection*{1.1 Course Information}

The lecture is part of CSE400: Fundamentals of Probability in Computing.

The lecture covers:

Introduction to probability theory

Joint probability

Conditional probability

Engineering applications

Logical dependency: \\
The lecture establishes foundational probability theory before introducing joint and conditional probability. 

L4

\subsection*{1.2 Why Learn Probability (As stated in lecture)}

Reasons given:

Daily life conversations

Engineering applications

Engineering applications listed:

Speech recognition

System radar system

Communication network

Logical reasoning presented: \\
These systems involve uncertainty and require reasoning about possible outcomes. \\
Therefore, probability is needed to quantify likelihood and support decisions. 

L4

\section*{2. Engineering Applications}

\subsection*{2.1 Speech Recognition System}

The lecture shows:

Vocabulary: Hello, Yes, No, Bye

Templates for different speakers

Logical dependency presented:

A spoken input may correspond to multiple templates.

The system must determine which word was spoken.

This requires evaluating likelihood of each possibility.

Probability provides this evaluation. 

L4

\subsection*{2.2 Radar System}

The lecture shows:

Radar sending signals

Signal reflections received

Logical dependency presented:

Reflected signals vary.

Detection must determine likelihood of object presence.

Probability is used to measure likelihood. 

L4

\subsection*{2.3 Communication Network}

The lecture shows:

Source and destination

Multiple communication paths

Logical dependency presented:

Multiple routes exist between nodes.

Transmission success is uncertain.

Probability models reliability and outcomes. 

L4

\section*{3. Structure of Lecture Content}

The lecture is organized into:

\subsection*{3.1 Introduction to Probability Theory}

Includes:

Experiments, sample space, and events

Axioms of probability

Corollaries and propositions

Assigning probability

\subsection*{3.2 Joint Probability}

Includes:

Motivation, notation, and concepts

Example: Card deck

Example: Costume party

\subsection*{3.3 Conditional Probability}

Includes:

Motivation, notation, and concepts

Example: Cards without replacement

Example: Game of poker

Example: Missing key

Logical dependency: \\
Joint and conditional probability require prior understanding of:

Experiment

Outcome

Event

Sample space

Probability axioms 

L4

\section*{4. Fundamental Definitions}

\subsection*{4.1 Experiment (E)}

Definition: \\
An experiment is a procedure we perform that produces some result.

Example given: \\
Tossing a coin five times (E₅)

Logical dependency:

An experiment produces results.

These results are outcomes.

Probability theory begins with experiments. 

L4

\subsection*{4.2 Outcome (ξ)}

Definition: \\
An outcome is a possible result of an experiment.

Example: \\
One outcome of E₅ is: \\
ξ₁ = HHTHT

Logical dependency:

Outcomes are generated by experiments.

Events are defined using outcomes. 

L4

\subsection*{4.3 Event}

Definition: \\
An event is a certain set of outcomes of an experiment.

Example: \\
Event C in experiment E₅: \\
C = \{all outcomes consisting of an even number of heads\}

Logical dependency:

Events are subsets of outcomes.

Probability is assigned to events. 

L4

\section*{5. Sample Space}

\subsection*{5.1 Definition of Sample Space (S)}

Definition: \\
Sample space is the collection or set of all possible distinct outcomes of an experiment.

Logical dependency:

All possible outcomes must be included.

Probability assignments require complete outcome set. 

L4

\subsection*{5.2 Properties of Sample Space}

Mutually Exclusive

Only one outcome can occur at a time.

Example: \\
You can get heads or tails, but not both.

Collectively Exhaustive

All possible outcomes must be included.

Example: \\
You cannot get anything other than heads or tails.

Logical dependency:

Outcomes must not overlap.

All possibilities must be covered.

This ensures valid probability assignment. 

L4

\subsection*{5.3 Nature of Sample Space}

Sample space is:

Discrete

Countably infinite

Continuous

Logical dependency: \\
The type of sample space determines how probability is represented. 

L4

\subsection*{5.4 Examples of Sample Spaces}

Examples listed in lecture:

Flipping a fair coin once

Cubical die with numbered faces rolled

Rolling two dice

Flip coin until a tails occurs

Random number generator with interval [0,1]

Logical purpose: \\
These illustrate types of sample spaces. 

L4

\section*{6. Probability}

\subsection*{6.1 Definition of Probability}

Probability is:

A measure of likelihood of various events, OR

A function of an event producing a numerical quantity measuring likelihood.

Logical dependency:

Events must have measurable likelihood.

Probability function provides numeric measure. 

L4

\section*{7. Axioms of Probability}

Axioms: \\
Statements taken to be self-evident and requiring no proof.

Logical role: \\
All results in probability follow from these axioms. 

L4

\subsection*{Axiom 1: Range of Probability}

For any event A: \\
$0 \leq Pr(A) \leq 1$

Logical dependency: \\
Probability is a measure of likelihood and must lie within this range. 

L4

\subsection*{Axiom 2: Probability of Sample Space}

If S is the sample space: \\
$Pr(S) = 1$

Logical dependency:

Sample space contains all outcomes.

One outcome must occur.

Total probability equals 1. 

L4

\subsection*{Axiom 3: Additivity for Mutually Exclusive Events}

If $A \cap B = \emptyset$: \\
$Pr(A \cup B) = Pr(A) + Pr(B)$

Logical dependency: \\
Mutually exclusive events cannot occur together; therefore probabilities add. 

L4

\subsection*{Infinite Additivity}

For mutually exclusive sets $A_i$: \\
$A_i \cap A_j = \emptyset$ for $i \neq j$

$Pr\left(\bigcup A_i\right) = \sum Pr(A_i)$

Logical dependency: \\
Extends additivity to infinite mutually exclusive events. 

L4

\section*{8. Corollary from Probability Axioms}

\subsection*{Corollary 2.1}

For M finite mutually exclusive events $A_i$: \\
$Pr\left(\bigcup A_i\right) = \sum Pr(A_i)$

Logical dependency:

Derived from Axiom 3.

Equivalent to Axiom 3 when sample space is finite.

Axiom 3 general form required for infinite sample space. 

L4

\section*{9. Propositions from Probability Axioms}

\subsection*{Proposition 2.1: Complement Rule}

$Pr(A^c) = 1 - Pr(A)$

Logical dependency:

Event and complement together form sample space.

From $Pr(S)=1$.

Remaining probability equals complement. 

L4

\subsection*{Proposition 2.2: Monotonicity}

If $A \subset B$: \\
$Pr(A) \leq Pr(B)$

Logical dependency:

Event B contains all outcomes of A.

Therefore probability of B cannot be less than probability of A. 

L4

\section*{10. Logical Dependency Chain of Entire Lecture}

The lecture builds in the following order:

Engineering problems motivate probability.

Probability theory begins with experiments.

Experiments produce outcomes.

Outcomes form events.

All outcomes form sample space.

Probability assigns likelihood to events.

Axioms define valid probability assignment.

Corollaries derived from axioms.

Propositions derived from axioms.

These foundations are required for:

Joint probability

Conditional probability

\end{document}
